\begin{lemma}[Countable hereditary-node closure]
Let $I$ be a nonempty countable type in the universe of $Q$, and let
$f:I\to\mathsf{HerCtblPower}(Q)$.  Then the full-hierarchy node
\[
\mathsf{node}\bigl(I,(f(i))_{i\in I}\bigr)
\]
is hereditarily countable.  This is a witness theorem about a node of the
full, unrestricted hierarchy; it does not replace $I$ by a countable index
in the definition of $\mathsf{PowerQ}$.
\end{lemma}
\begin{proof}
By countability and nonemptiness of $I$, choose a surjection
$e:\mathbb N\twoheadrightarrow I$, using the standard countable-surjection
fact.  The code objects used below are the objects defined in
$\noderef{HerCtblCode}$, and their erasures are those of
$\noderef{HerCtblCodeErase}$.  By the subtype definition
$\noderef{HerCtblPower}$, each $f(i)$ has a hereditary-countability witness;
then $\noderef{HereditarilyCountable}$ provides a code $c_i$ such that
$\mathsf{PowerPresentationEq}(f(i),\mathsf{erase}(c_i))$.  Choose one such
code for each $i$ and form the code $c$ whose $n$th child is $c_{e(n)}$.

The erasure of $c$ is a node indexed by $\mathsf{ULift}(\mathbb N)$.
To compare it with the original $I$-indexed node, apply the two-node
clause of $\noderef{PowerPresentationEq}$.  Given $i\in I$, surjectivity
gives $n$ with $e(n)=i$; use $\mathsf{ULift.up}(n)$ and the chosen
equivalence for $c_i$.  Conversely, a lifted $n$ is matched by the index
$e(n)$, using the same chosen equivalence in the reverse direction, supplied
by symmetry in $\noderef{PowerPresentationEqEquivalence}$.  These witnesses
establish presentation equivalence between the full node and
$\mathsf{erase}(c)$, which is exactly the defining condition
\noderef{HereditarilyCountable}.
\end{proof}
